\subsection{A reduced power whose nontrivial quotients all have continuum size}

\begin{claim}{A.24}{Proved}{Reduced powers need not admit a smaller detecting quotient}
Let $Q$ be a countably infinite torsion-free group in which all
nonidentity elements are conjugate, as constructed in A.21. Put
\[
 R=Q^{\mathbb N}/Q^{(\mathbb N)}.
 \tag{A.24.1}
\]
Every nontrivial quotient of $R$ has cardinality exactly
$2^{\aleph_0}$. More precisely, every such quotient admits a
specified epimorphism onto a nonprincipal ultrapower of $Q$.
Consequently every homomorphism from $R$ to a group of cardinality
strictly below $2^{\aleph_0}$ is trivial.

The diagonal map $Q\hookrightarrow R$ is injective, $R$ is
torsion-free, and for every positive integer $m$ there is an
explicit embedding $\mathbb Z^m\hookrightarrow R$. In particular
$R$ has no finite-dimensional classifying CW complex. Thus this
example gives a precise obstruction to reducing these particular
reduced-power targets to countable quotients; it does not assert
the same quotient property for every possible choice of $E$ in
A.19.
\end{claim}

\begin{proof}
Write $P=Q^{\mathbb N}$ and, for $g=(g_n)\in P$, put
$\operatorname{supp}(g)=\{n:g_n\ne1\}$. We first determine every
normal subgroup of $P$ by these supports.

For $g,h\in P$ with
$\operatorname{supp}(h)\subseteq\operatorname{supp}(g)$, the
element $h$ is a product of two conjugates of $g$. To see this,
at a coordinate where $g_n\ne1$, choose
$a_n\in Q\setminus\{1,h_n\}$ and put $b_n=a_n^{-1}h_n$.
Both $a_n$ and $b_n$ are nonidentity. They are therefore conjugate
to $g_n$, so choose $u_n,v_n\in Q$ with
\[
 a_n=u_ng_nu_n^{-1},\qquad b_n=v_ng_nv_n^{-1}.
\]
At a coordinate where $g_n=1$, also $h_n=1$, and take
$u_n=v_n=1$. These choices give the actual equality
\[
 h=ugu^{-1}vgv^{-1}\quad\text{in }P.
 \tag{A.24.2}
\]
It follows that a normal subgroup containing one element of
support $S$ contains every element supported in $S$.

For $K\triangleleft P$, let
\[
 \mathcal I_K=\{S\subseteq\mathbb N:
       \text{some element of }K\text{ has support }S\}.
\]
Every subset of $\mathbb N$ is the support of a sequence, by
using a fixed nonidentity element of $Q$ on that subset.
Formula (A.24.2) shows that $\mathcal I_K$ is closed under
taking subsets and that
\[
 K=\{g\in P:\operatorname{supp}(g)\in\mathcal I_K\}.
 \tag{A.24.3}
\]
It is also closed under finite unions. Indeed, for
$S,T\in\mathcal I_K$, any sequence supported in $S\cup T$
is a product of its restriction to $S$ and its restriction to
$T\setminus S$, and both factors belong to $K$ by (A.24.2).
In particular a sequence with support exactly $S\cup T$ belongs
to $K$. Thus $\mathcal I_K$ is an ideal of subsets. Conversely,
for any such ideal $\mathcal I$, the set in the right side of
(A.24.3) is a normal subgroup: inverses and conjugation preserve
support, and a product's support is contained in the union of
the two supports. These two constructions are inverse. In
particular $K$ is proper exactly when $\mathcal I_K$ is proper.

Let $R\twoheadrightarrow V$ be a nontrivial quotient and let
$K$ be the kernel of the composite $P\to R\to V$. Then $K$
is proper and contains every finite-support sequence, so
$\mathcal I_K$ is a proper ideal containing all finite subsets.
Extend it to a maximal proper ideal $\mathcal M$ of
$\mathcal P(\mathbb N)$. Its complementary family
$\mathcal U=\mathcal P(\mathbb N)\setminus\mathcal M$ is an
ultrafilter. It contains no finite set, since all finite sets
belong to $\mathcal M$, and is therefore nonprincipal.

The kernel of the canonical map
$P\to Q^{\mathbb N}/\mathcal U$ is
\[
 \{g:\{n:g_n=1\}\in\mathcal U\}
   =\{g:\operatorname{supp}(g)\in\mathcal M\}.
\]
It contains $K$. Hence the explicit formula
\[
 V=P/K\longrightarrow Q^{\mathbb N}/\mathcal U,
       \qquad gK\longmapsto[g]_{\mathcal U}
 \tag{A.24.4}
\]
is a well-defined epimorphism.

This ultrapower has cardinality exactly $2^{\aleph_0}$. The
upper bound follows from countability of $Q$. For the lower
bound, choose an injection $c$ from the finite binary words
into $Q$. To an infinite binary sequence $s$ associate the
sequence $n\mapsto c(s|_n)$. For distinct $s,t$, their prefixes
differ for every sufficiently large $n$, so the equality set
of these two $Q$-sequences is finite and does not belong to
$\mathcal U$. Thus they give distinct ultrapower elements.
This embeds the set of infinite binary sequences. Formula
(A.24.4) gives $|V|\geq2^{\aleph_0}$, while
$|V|\leq|R|\leq|P|=2^{\aleph_0}$ gives equality. Applying
this conclusion to the image of any homomorphism proves
the asserted vanishing for smaller targets.

A nonidentity constant sequence has infinite support, proving
that the diagonal $Q\to R$ is injective. If a positive power
of a sequence represents the identity in $R$, its coordinate
powers are the identity for all but finitely many coordinates.
Torsion-freeness of $Q$ makes the original sequence finite
supported, proving torsion-freeness of $R$.

For the dimension assertion, fix $a\in Q\setminus\{1\}$
and partition $\mathbb N$ into $m$ infinite subsets
$S_1,\ldots,S_m$. Define
\[
 \mathbb Z^m\longrightarrow R,
 \qquad (k_1,\ldots,k_m)\longmapsto
 \bigl[(a^{k_j})_{n\in S_j,\ 1\leq j\leq m}\bigr].
 \tag{A.24.5}
\]
Multiplication is coordinatewise, so this is a homomorphism.
If any $k_j\ne0$, torsion-freeness of $Q$ makes its support
contain the infinite set $S_j$, so its class in $R$ is not
the identity. Hence (A.24.5) is injective. A $d$-dimensional
classifying CW complex for $R$ would have a $d$-dimensional
cover classifying each such subgroup. Taking $m>d$ contradicts
$H^m(\mathbb Z^m;\mathbb Z)\cong\mathbb Z$, as computed
by the $m$-torus. Thus no finite $d$ is possible.
\end{proof}

